\documentclass{ximera}

\input{../../../preamble.tex}


\definecolor{fvdnavy}{HTML}{071D43}
\definecolor{fvdcyan}{HTML}{20BFC6}
\definecolor{fvdcyanlight}{HTML}{EFFCFB}
\definecolor{fvdmagenta}{HTML}{F10D69}
\definecolor{fvdmagentalight}{HTML}{FFF1F7}
\definecolor{fvdtext}{HTML}{163044}





%=================================================
% Pedagogische omgevingen
% PDF: tcolorbox
% HTML: learning-box
%=================================================

\newenvironment{voorkennis}
{%
    \ifdefined\HCode
        \HCode{%
<section class="learning-box learning-box--prior">
<div class="learning-box__header">
<span class="learning-box__icon" aria-hidden="true"></span>
<span class="learning-box__title">Voorkennis</span>
</div>
<div class="learning-box__content">
        }%
        \begin{itemize}
    \else
       \begin{tcolorbox}[
    enhanced,
    breakable,
    colback=fvdcyanlight,
    colframe=fvdcyan,
    coltext=fvdtext,
    boxrule=0.5pt,
    leftrule=4pt,
    arc=4mm,
    outer arc=4mm,
    left=7mm,
    right=6mm,
    top=5mm,
    bottom=5mm,
    before skip=6mm,
    after skip=6mm,
    title=Voorkennis,
    coltitle=fvdcyan,
    fonttitle=\bfseries\Large,
    attach title to upper={\par\smallskip},
    boxed title style={
        empty,
        left=0mm,
        right=0mm,
        top=0mm,
        bottom=0mm
    },
    overlay={
        \draw[
            fvdcyan,
            line width=0.5pt
        ]
        ([xshift=42mm,yshift=-13mm]frame.north west)
        --
        ([xshift=-7mm,yshift=-13mm]frame.north east);
    }
]
        \begin{itemize}
    \fi
}
{%
    \end{itemize}
    \ifdefined\HCode
        \HCode{%
</div>
</section>
        }%
    \else
        \end{tcolorbox}
    \fi
}


\newenvironment{leerdoelen}
{%
    \ifdefined\HCode
        \HCode{%
<section class="learning-box learning-box--goals">
<div class="learning-box__header">
<span class="learning-box__icon" aria-hidden="true"></span>
<span class="learning-box__title">Leerdoelen</span>
</div>
<div class="learning-box__content">
        }%
        \begin{itemize}
    \else
\begin{tcolorbox}[
    enhanced,
    breakable,
    colback=fvdmagentalight,
    colframe=fvdmagenta,
    coltext=fvdtext,
    boxrule=0.5pt,
    leftrule=4pt,
    arc=4mm,
    outer arc=4mm,
    left=7mm,
    right=6mm,
    top=5mm,
    bottom=5mm,
    before skip=6mm,
    after skip=6mm,
    title=Leerdoelen,
    coltitle=fvdmagenta,
    fonttitle=\bfseries\Large,
    attach title to upper={\par\smallskip},
    boxed title style={
        empty,
        left=0mm,
        right=0mm,
        top=0mm,
        bottom=0mm
    },
    overlay={
        \draw[
            fvdmagenta,
            line width=0.5pt
        ]
        ([xshift=42mm,yshift=-13mm]frame.north west)
        --
        ([xshift=-7mm,yshift=-13mm]frame.north east);
    }
]
        \begin{itemize}
    \fi
}
{%
    \end{itemize}
    \ifdefined\HCode
        \HCode{%
</div>
</section>
        }%
    \else
        \end{tcolorbox}
    \fi
}


\addPrintStyle{../../..}

\title{Oefeningen — Verloop van functies}
\author{Ingmar Herreman}



\begin{document}

% FVD_CHAPTER_DATA_START
% Automatisch gegenereerd uit index.tex.
% Niet handmatig aanpassen.
\fvdchapterdata{1}{Veranderingen in beeld}{3}
% FVD_CHAPTER_DATA_END














































































\maketitle

\FVDbeginExercises

% ============================================================
% BASIS
% ============================================================




% ------------------------------------------------------------
% OEFENING 1
% ------------------------------------------------------------

\begin{oefening}[Van het teken van $f'$ naar het verloop]{\oefB\ESS}

Van een functie $f$ is de volgende informatie over $f'$ gegeven:

\[
\begin{array}{c|ccccc}
x      & -\infty &   & -2 &   & +\infty \\ \hline
f'(x)  &         & + & 0  & - &
\end{array}
\]

\begin{question}
Op welk interval is $f$ stijgend?
\end{question}

\begin{question}
Op welk interval is $f$ dalend?
\end{question}

\begin{question}
Wat kun je besluiten over $x=-2$?
\end{question}

\begin{question}
Gaat het om een lokaal maximum of een lokaal minimum?
Motiveer met de tekenverandering van $f'$.
\end{question}

\begin{oplossing}

Uit de tekentabel volgt:
\[
f'(x)>0 \quad \text{voor } x<-2,
\]
dus $f$ is stijgend op
\[
]-\infty,-2[.
\]

Verder:
\[
f'(x)<0 \quad \text{voor } x>-2,
\]
dus $f$ is dalend op
\[
]-2,+\infty[.
\]

Bij $x=-2$ geldt
\[
f'(-2)=0.
\]
De afgeleide verandert daar van positief naar negatief. De functie gaat
dus over van stijgen naar dalen.

Bijgevolg heeft $f$ bij $x=-2$ een lokaal maximum.

De functiewaarde van dit maximum kunnen we uit de gegeven informatie niet
bepalen.

\end{oplossing}

\end{oefening}


% ------------------------------------------------------------
% OEFENING 2
% ------------------------------------------------------------

\begin{oefening}[Minimum herkennen]{\oefB\ESS}

Voor een functie $g$ geldt:

\[
\begin{array}{c|ccccc}
x      & -\infty &   & 3 &   & +\infty \\ \hline
g'(x)  &         & - & 0 & + &
\end{array}
\]

\begin{question}
Waar is $g$ dalend?
\end{question}

\begin{question}
Waar is $g$ stijgend?
\end{question}

\begin{question}
Wat gebeurt er bij $x=3$?
\end{question}

\begin{question}
Kun je uit deze gegevens de functiewaarde van het extremum bepalen?
Leg uit.
\end{question}

\begin{oplossing}

Uit de tekentabel volgt:
\[
g'(x)<0 \quad \text{voor } x<3,
\]
dus $g$ is dalend op
\[
]-\infty,3[.
\]

Voor $x>3$ geldt
\[
g'(x)>0,
\]
dus $g$ is stijgend op
\[
]3,+\infty[.
\]

Bij $x=3$ verandert $g'$ van negatief naar positief. De functie gaat dus
over van dalen naar stijgen en heeft daar een lokaal minimum.

De functiewaarde van het minimum kunnen we niet bepalen, omdat $g(3)$ niet
gegeven is.

\end{oplossing}

\end{oefening}


% ------------------------------------------------------------
% OEFENING 3
% ------------------------------------------------------------

\begin{oefening}[Stationair is niet altijd extreem]{\oefB\ESS}

Van drie functies kennen we het tekenverloop van de afgeleide.

\[
\begin{array}{c|ccccc}
x      &      &   & a &   &      \\ \hline
f'(x)  &      & + & 0 & - &
\end{array}
\]

\[
\begin{array}{c|ccccc}
x      &      &   & b &   &      \\ \hline
g'(x)  &      & - & 0 & + &
\end{array}
\]

\[
\begin{array}{c|ccccc}
x      &      &   & c &   &      \\ \hline
h'(x)  &      & + & 0 & + &
\end{array}
\]

\begin{question}
Welke functie heeft een lokaal maximum?
\end{question}

\begin{question}
Welke functie heeft een lokaal minimum?
\end{question}

\begin{question}
Bij welke functie is er wel een stationair punt maar geen extremum?
\end{question}

\begin{question}
Waarom volstaat de voorwaarde
\[
f'(a)=0
\]
niet om te besluiten dat $f$ in $a$ een extremum heeft?
\end{question}

\begin{oplossing}

Bij $x=a$ verandert $f'$ van positief naar negatief. Daarom heeft $f$ daar
een lokaal maximum.

Bij $x=b$ verandert $g'$ van negatief naar positief. Daarom heeft $g$ daar
een lokaal minimum.

Bij $x=c$ geldt
\[
h'(c)=0,
\]
maar $h'$ is zowel links als rechts van $c$ positief. De functie blijft dus
stijgen. Er is wel een stationair punt, maar geen extremum.

De voorwaarde
\[
f'(a)=0
\]
zegt alleen dat de raaklijn horizontaal is. Voor een extremum moet bovendien
het teken van $f'$ links en rechts van $a$ veranderen.

\end{oplossing}

\end{oefening}


% ------------------------------------------------------------
% OEFENING 4
% ------------------------------------------------------------

\begin{oefening}[Stijgen en dalen onderzoeken]{\oefB\ESS}

Gegeven:

\[
f(x)=x^2-4x+1.
\]

\begin{question}
Bepaal $f'(x)$.
\end{question}

\begin{question}
Los op:
\[
f'(x)=0.
\]
\end{question}

\begin{question}
Stel een tekentabel op voor $f'$.
\end{question}

\begin{question}
Bepaal de intervallen waarop $f$ stijgt en daalt.
\end{question}

\begin{question}
Bepaal het lokale extremum van $f$.
\end{question}

\begin{oplossing}

Gegeven:
\[
f(x)=x^2-4x+1.
\]

De afgeleide is
\[
f'(x)=2x-4=2(x-2).
\]

Dus
\[
f'(x)=0
\quad\Longleftrightarrow\quad
x=2.
\]

Het tekenverloop is:
\[
\begin{array}{c|ccccc}
x & -\infty & & 2 & & +\infty\\ \hline
f'(x) & & - & 0 & + &
\end{array}
\]

Daarom is $f$ dalend op
\[
]-\infty,2[
\]
en stijgend op
\[
]2,+\infty[.
\]

Bij $x=2$ heeft $f$ een lokaal minimum. De functiewaarde is
\[
f(2)=4-8+1=-3.
\]

Het lokale minimum is dus
\[
\boxed{(2,-3)}.
\]

\end{oplossing}

\end{oefening}


% ------------------------------------------------------------
% OEFENING 5
% ------------------------------------------------------------

\begin{oefening}[Een derdegraadsfunctie onderzoeken]{\oefB\ESS}

Gegeven:

\[
f(x)=x^3-3x^2-9x+5.
\]

\begin{question}
Bepaal $f'(x)$.
\end{question}

\begin{question}
Factoriseer $f'(x)$.
\end{question}

\begin{question}
Bepaal de stationaire punten.
\end{question}

\begin{question}
Stel een tekentabel op voor $f'$.
\end{question}

\begin{question}
Bepaal waar $f$ stijgt en waar $f$ daalt.
\end{question}

\begin{question}
Bepaal de lokale extrema en geef hun coördinaten.
\end{question}

\begin{oplossing}

Gegeven:
\[
f(x)=x^3-3x^2-9x+5.
\]

De afgeleide is
\[
f'(x)=3x^2-6x-9.
\]

Factoriseren geeft
\[
f'(x)=3(x^2-2x-3)
=3(x+1)(x-3).
\]

De stationaire waarden zijn dus
\[
x=-1
\qquad\text{en}\qquad
x=3.
\]

Het tekenverloop:
\[
\begin{array}{c|ccccccc}
x & -\infty & & -1 & & 3 & & +\infty\\ \hline
f'(x) & & + & 0 & - & 0 & + &
\end{array}
\]

Daarom is $f$ stijgend op
\[
]-\infty,-1[
\quad\text{en}\quad
]3,+\infty[
\]
en dalend op
\[
]-1,3[.
\]

Bij $x=-1$ verandert $f'$ van positief naar negatief, dus daar ligt een
lokaal maximum:
\[
f(-1)=-1-3+9+5=10.
\]

Bij $x=3$ verandert $f'$ van negatief naar positief, dus daar ligt een
lokaal minimum:
\[
f(3)=27-27-27+5=-22.
\]

De lokale extrema zijn dus
\[
\boxed{(-1,10)}
\qquad\text{en}\qquad
\boxed{(3,-22)}.
\]

\end{oplossing}

\end{oefening}


% ------------------------------------------------------------
% OEFENING 6
% ------------------------------------------------------------

\begin{oefening}[Geen extremum ondanks $f'(a)=0$]{\oefB\ESS}

Gegeven:

\[
f(x)=x^3.
\]

\begin{question}
Bepaal $f'(x)$.
\end{question}

\begin{question}
Voor welke $x$ geldt $f'(x)=0$?
\end{question}

\begin{question}
Onderzoek het teken van $f'$ links en rechts van deze waarde.
\end{question}

\begin{question}
Heeft $f$ daar een lokaal extremum?
\end{question}

\begin{question}
Welke fout maakt een leerling die redeneert:
``$f'(0)=0$, dus $(0,0)$ is een extremum''?
\end{question}

\begin{oplossing}

Voor
\[
f(x)=x^3
\]
geldt
\[
f'(x)=3x^2.
\]

Dus
\[
f'(x)=0
\quad\Longleftrightarrow\quad
x=0.
\]

Voor zowel $x<0$ als $x>0$ geldt
\[
f'(x)>0.
\]

De functie is dus aan beide kanten van $0$ stijgend. Bij $x=0$ is er geen
tekenverandering van $f'$ en bijgevolg geen lokaal extremum.

De leerling maakt de fout om
\[
f'(0)=0
\]
als voldoende voorwaarde voor een extremum te beschouwen. Het is slechts
een kandidaatvoorwaarde; het tekenverloop rond het punt moet nog onderzocht
worden.

\end{oplossing}

\end{oefening}


% ------------------------------------------------------------
% OEFENING 7
% ------------------------------------------------------------

\begin{oefening}[Van $f''$ naar kromming]{\oefB\ESS}

Van een functie $f$ is gegeven:

\[
\begin{array}{c|ccccc}
x       & -\infty &   & 2 &   & +\infty \\ \hline
f''(x)  &         & - & 0 & + &
\end{array}
\]

\begin{question}
Beschrijf de kromming van de grafiek voor $x<2$.
\end{question}

\begin{question}
Beschrijf de kromming van de grafiek voor $x>2$.
\end{question}

\begin{question}
Wat kun je besluiten bij $x=2$?
\end{question}

\begin{question}
Welke rol speelt de tekenverandering van $f''$ in je besluit?
\end{question}

\begin{oplossing}

Voor $x<2$ geldt
\[
f''(x)<0.
\]
De grafiek is daar bol (concaaf) en $f'$ is dalend.

Voor $x>2$ geldt
\[
f''(x)>0.
\]
De grafiek is daar hol (convex) en $f'$ is stijgend.

Bij $x=2$ verandert $f''$ van teken. De kromming van de grafiek verandert
dus. Bijgevolg heeft de grafiek bij abscis $x=2$ een buigpunt.

De $y$-coördinaat kan niet worden bepaald zonder de functiewaarde $f(2)$.

\end{oplossing}

\end{oefening}


% ------------------------------------------------------------
% OEFENING 8
% ------------------------------------------------------------

\begin{oefening}[Buigpunt onderzoeken]{\oefB\ESS}

Gegeven:

\[
f(x)=x^3-6x^2+4x+3.
\]

\begin{question}
Bepaal $f'(x)$.
\end{question}

\begin{question}
Bepaal $f''(x)$.
\end{question}

\begin{question}
Los op:
\[
f''(x)=0.
\]
\end{question}

\begin{question}
Onderzoek het teken van $f''$.
\end{question}

\begin{question}
Heeft de grafiek een buigpunt?
Motiveer.
\end{question}

\begin{question}
Bepaal de coördinaten van het buigpunt.
\end{question}

\begin{oplossing}

Gegeven:
\[
f(x)=x^3-6x^2+4x+3.
\]

De eerste afgeleide is
\[
f'(x)=3x^2-12x+4.
\]

De tweede afgeleide is
\[
f''(x)=6x-12=6(x-2).
\]

Dus
\[
f''(x)=0
\quad\Longleftrightarrow\quad
x=2.
\]

Voor $x<2$ geldt
\[
f''(x)<0,
\]
en voor $x>2$ geldt
\[
f''(x)>0.
\]

De tweede afgeleide verandert dus van teken. De grafiek heeft bijgevolg
een buigpunt bij $x=2$.

De functiewaarde is
\[
f(2)=8-24+8+3=-5.
\]

Het buigpunt is
\[
\boxed{(2,-5)}.
\]

\end{oplossing}

\end{oefening}


% ------------------------------------------------------------
% OEFENING 9
% ------------------------------------------------------------

\begin{oefening}[$f''(a)=0$ volstaat niet]{\oefB\ESS}

Gegeven:

\[
f(x)=x^4.
\]

\begin{question}
Bepaal $f''(x)$.
\end{question}

\begin{question}
Toon aan dat
\[
f''(0)=0.
\]
\end{question}

\begin{question}
Onderzoek het teken van $f''$ links en rechts van $0$.
\end{question}

\begin{question}
Is $(0,0)$ een buigpunt?
\end{question}

\begin{question}
Formuleer op basis hiervan waarom $f''(a)=0$ alleen
niet voldoende is om een buigpunt te besluiten.
\end{question}

\begin{oplossing}

Voor
\[
f(x)=x^4
\]
geldt
\[
f'(x)=4x^3
\]
en
\[
f''(x)=12x^2.
\]

Daarom
\[
f''(0)=0.
\]

Maar voor zowel $x<0$ als $x>0$ geldt
\[
f''(x)>0.
\]

De tweede afgeleide verandert dus niet van teken en de kromming verandert
niet. Bijgevolg is
\[
(0,0)
\]
geen buigpunt.

De voorwaarde
\[
f''(a)=0
\]
is dus op zichzelf niet voldoende. We moeten bijkomend onderzoeken of de
kromming, of equivalent het teken van $f''$, werkelijk verandert.

\end{oplossing}

\end{oefening}


% ------------------------------------------------------------
% OEFENING 10
% ------------------------------------------------------------

\begin{oefening}[Volledig verlooponderzoek]{\oefB\ESS}

Gegeven:

\[
f(x)=x^3-3x.
\]

\begin{question}
Bepaal $f'(x)$ en $f''(x)$.
\end{question}

\begin{question}
Bepaal de stationaire punten.
\end{question}

\begin{question}
Stel een tekentabel op voor $f'$.
\end{question}

\begin{question}
Bepaal de stijg- en daalintervallen.
\end{question}

\begin{question}
Bepaal de lokale extrema.
\end{question}

\begin{question}
Onderzoek het teken van $f''$.
\end{question}

\begin{question}
Bepaal het buigpunt.
\end{question}

\begin{question}
Maak op basis van je onderzoek een kwalitatieve schets van de grafiek.
\end{question}

\begin{oplossing}

Gegeven:
\[
f(x)=x^3-3x.
\]

We vinden:
\[
f'(x)=3x^2-3=3(x-1)(x+1),
\]
en
\[
f''(x)=6x.
\]

De stationaire waarden zijn
\[
x=-1
\qquad\text{en}\qquad
x=1.
\]

Het tekenverloop van $f'$ is:
\[
\begin{array}{c|ccccccc}
x & -\infty & & -1 & & 1 & & +\infty\\ \hline
f'(x) & & + & 0 & - & 0 & + &
\end{array}
\]

Dus $f$ stijgt op
\[
]-\infty,-1[
\quad\text{en}\quad
]1,+\infty[
\]
en daalt op
\[
]-1,1[.
\]

Bij $x=-1$ ligt een lokaal maximum:
\[
f(-1)=-1+3=2.
\]
Dus
\[
\boxed{(-1,2)}.
\]

Bij $x=1$ ligt een lokaal minimum:
\[
f(1)=1-3=-2.
\]
Dus
\[
\boxed{(1,-2)}.
\]

Voor de kromming:
\[
f''(x)<0 \quad \text{voor }x<0,
\]
en
\[
f''(x)>0 \quad \text{voor }x>0.
\]

De kromming verandert bij $x=0$. Omdat
\[
f(0)=0,
\]
is het buigpunt
\[
\boxed{(0,0)}.
\]

De kwalitatieve grafiek is een derdegraadsgrafiek die stijgt tot
$(-1,2)$, vervolgens daalt via het buigpunt $(0,0)$ tot $(1,-2)$ en daarna
opnieuw stijgt.

\end{oplossing}

\end{oefening}


% ============================================================
% VERDIEPING
% ============================================================




% ------------------------------------------------------------
% OEFENING 11
% ------------------------------------------------------------

\begin{oefening}[Van grafiek naar afgeleide]{\oefU\ESS}

Hieronder zie je de grafiek van een functie $f$.

\begin{center}
\begin{tikzpicture}[x=1cm,y=0.65cm]

  \draw[->,draw=black!45] (-4,0)--(4.3,0)
    node[right] {$x$};
  \draw[->,draw=black!45] (0,-4)--(0,4.5)
    node[above] {$y$};

  \draw[
    draw=fvdnavy,
    very thick,
    smooth,
    domain=-3:3,
    samples=150
  ]
    plot (\x,{0.35*(\x)^3-1.4*\x});

\end{tikzpicture}
\end{center}

Gebruik uitsluitend de grafiek.

\begin{question}
Waar is $f$ stijgend?
\end{question}

\begin{question}
Waar is $f$ dalend?
\end{question}

\begin{question}
Voor welke waarden van $x$ verwacht je
\[
f'(x)=0?
\]
\end{question}

\begin{question}
Geef aan waar $f'(x)>0$ en waar $f'(x)<0$.
\end{question}

\begin{question}
Schets kwalitatief een mogelijke grafiek van $f'$.
\end{question}

\begin{oplossing}

Uit de grafiek zien we een lokaal maximum rond
\[
x\approx -1.15
\]
en een lokaal minimum rond
\[
x\approx 1.15.
\]

Daarom is $f$ stijgend voor ongeveer
\[
x<-1.15
\quad\text{en}\quad
x>1.15,
\]
en dalend voor ongeveer
\[
-1.15<x<1.15.
\]

Bij de twee extrema verwachten we
\[
f'(x)=0.
\]

Dus kwalitatief:
\[
f'(x)>0
\]
buiten deze twee waarden en
\[
f'(x)<0
\]
ertussen.

Een mogelijke grafiek van $f'$ is bijgevolg een naar boven geopende
parabool met nulpunten rond $x=-1.15$ en $x=1.15$.

\end{oplossing}

\end{oefening}


% ------------------------------------------------------------
% OEFENING 12
% ------------------------------------------------------------

\begin{oefening}[Van $f'$ naar $f$]{\oefU\ESS}

Hieronder zie je de grafiek van $f'$.

\begin{center}
\begin{tikzpicture}[x=1cm,y=0.7cm]

  \draw[->,draw=black!45] (-4,0)--(4.3,0)
    node[right] {$x$};
  \draw[->,draw=black!45] (0,-2)--(0,4.5)
    node[above] {$f'(x)$};

  \draw[
    draw=fvdnavy,
    very thick,
    smooth,
    domain=-3.2:3.2,
    samples=150
  ]
    plot (\x,{0.6*(\x+2)*(\x-1)});

  \fill (-2,0) circle (2pt);
  \fill (1,0) circle (2pt);

\end{tikzpicture}
\end{center}

\begin{question}
Voor welke waarden van $x$ is $f'(x)=0$?
\end{question}

\begin{question}
Waar is $f$ stijgend?
\end{question}

\begin{question}
Waar is $f$ dalend?
\end{question}

\begin{question}
Bij welke waarde van $x$ heeft $f$ een lokaal maximum?
\end{question}

\begin{question}
Bij welke waarde van $x$ heeft $f$ een lokaal minimum?
\end{question}

\begin{question}
De grafiek van $f'$ heeft zelf een minimum.
Wat vertelt dat over $f''$?
\end{question}

\begin{oplossing}

De grafiek van $f'$ snijdt de $x$-as bij
\[
x=-2
\qquad\text{en}\qquad
x=1.
\]

Dus
\[
f'(-2)=0
\qquad\text{en}\qquad
f'(1)=0.
\]

Omdat de parabool naar boven opent, geldt
\[
f'(x)>0
\]
voor
\[
x<-2
\quad\text{en}\quad
x>1,
\]
en
\[
f'(x)<0
\]
voor
\[
-2<x<1.
\]

Daarom is $f$ stijgend op
\[
]-\infty,-2[
\quad\text{en}\quad
]1,+\infty[
\]
en dalend op
\[
]-2,1[.
\]

Bij $x=-2$ verandert $f'$ van positief naar negatief, dus heeft $f$ daar
een lokaal maximum.

Bij $x=1$ verandert $f'$ van negatief naar positief, dus heeft $f$ daar
een lokaal minimum.

De grafiek van $f'$ heeft zelf een minimum halverwege haar nulpunten:
\[
x=\frac{-2+1}{2}=-\frac12.
\]

Daar is de helling van $f'$ nul, dus
\[
f''\left(-\frac12\right)=0.
\]
Omdat $f'$ daar van dalen naar stijgen gaat, verandert $f''$ van teken.
Voor $f$ komt dit overeen met een verandering van kromming.

\end{oplossing}

\end{oefening}


% ------------------------------------------------------------
% OEFENING 13
% ------------------------------------------------------------

\begin{oefening}[Drie functies verbinden]{\oefU\ESS}

Beschouw een functie $f$ en haar afgeleiden.

Op een interval $I$ geldt:

\[
f'(x)>0
\qquad\text{en}\qquad
f''(x)<0.
\]

\begin{question}
Wat weet je over het stijgen of dalen van $f$?
\end{question}

\begin{question}
Wat weet je over de kromming van de grafiek van $f$?
\end{question}

\begin{question}
Neemt $f'$ op dit interval toe of af?
\end{question}

\begin{question}
Beschrijf in woorden hoe $f$ op dit interval verandert.
\end{question}

\begin{oplossing}

Omdat
\[
f'(x)>0,
\]
is $f$ stijgend op het interval $I$.

Omdat
\[
f''(x)<0,
\]
is de grafiek bol (concaaf) en is $f'$ dalend.

De functie stijgt dus nog steeds, maar de richtingscoëfficiënten van de
raaklijnen worden kleiner.

Met andere woorden: $f$ stijgt steeds minder snel.

\end{oplossing}

\end{oefening}


% ------------------------------------------------------------
% OEFENING 14
% ------------------------------------------------------------

\begin{oefening}[Vier mogelijke situaties]{\oefU\ESS}

Beschrijf telkens zo nauwkeurig mogelijk het gedrag van $f$.

\begin{question}
\[
f'(x)>0,\qquad f''(x)>0.
\]
\end{question}

\begin{question}
\[
f'(x)>0,\qquad f''(x)<0.
\]
\end{question}

\begin{question}
\[
f'(x)<0,\qquad f''(x)>0.
\]
\end{question}

\begin{question}
\[
f'(x)<0,\qquad f''(x)<0.
\]
\end{question}

Geef bij elke situatie zowel informatie over het verloop van $f$
als over de verandering van de hellingen.

\begin{oplossing}

\textbf{1.}
\[
f'(x)>0,\qquad f''(x)>0.
\]

$f$ stijgt en $f'$ neemt toe. De hellingen worden dus steeds groter:
$f$ stijgt steeds sneller. De grafiek is hol (convex).

\medskip

\textbf{2.}
\[
f'(x)>0,\qquad f''(x)<0.
\]

$f$ stijgt, maar $f'$ neemt af. De hellingen worden kleiner:
$f$ stijgt steeds minder snel. De grafiek is bol (concaaf).

\medskip

\textbf{3.}
\[
f'(x)<0,\qquad f''(x)>0.
\]

$f$ daalt, maar $f'$ neemt toe. De negatieve hellingen worden minder
negatief: $f$ daalt steeds minder snel. De grafiek is hol (convex).

\medskip

\textbf{4.}
\[
f'(x)<0,\qquad f''(x)<0.
\]

$f$ daalt en $f'$ neemt verder af. De hellingen worden steeds negatiever:
$f$ daalt steeds sneller. De grafiek is bol (concaaf).

\end{oplossing}

\end{oefening}


% ------------------------------------------------------------
% OEFENING 15
% ------------------------------------------------------------

\begin{oefening}[Een volledige analyse]{\oefU\ESS}

Gegeven:

\[
f(x)=x^4-4x^2.
\]

\begin{question}
Bepaal $f'(x)$ en factoriseer.
\end{question}

\begin{question}
Bepaal alle stationaire punten.
\end{question}

\begin{question}
Onderzoek het teken van $f'$ en bepaal de stijg- en daalintervallen.
\end{question}

\begin{question}
Bepaal alle lokale extrema.
\end{question}

\begin{question}
Bepaal $f''(x)$.
\end{question}

\begin{question}
Onderzoek de kromming.
\end{question}

\begin{question}
Bepaal eventuele buigpunten.
\end{question}

\begin{question}
Maak een kwalitatieve grafiek en controleer of alle gevonden
eigenschappen zichtbaar zijn.
\end{question}

\begin{oplossing}

Gegeven:
\[
f(x)=x^4-4x^2.
\]

De eerste afgeleide is
\[
f'(x)=4x^3-8x
=4x(x^2-2)
=4x(x-\sqrt2)(x+\sqrt2).
\]

De stationaire waarden zijn
\[
x=-\sqrt2,\qquad x=0,\qquad x=\sqrt2.
\]

Het tekenverloop:
\[
\begin{array}{c|ccccccccc}
x & -\infty & & -\sqrt2 & & 0 & & \sqrt2 & & +\infty\\ \hline
f'(x) & & - & 0 & + & 0 & - & 0 & + &
\end{array}
\]

Dus $f$ daalt op
\[
]-\infty,-\sqrt2[
\quad\text{en}\quad
]0,\sqrt2[
\]
en stijgt op
\[
]-\sqrt2,0[
\quad\text{en}\quad
]\sqrt2,+\infty[.
\]

Bij $x=\pm\sqrt2$ liggen lokale minima:
\[
f(\pm\sqrt2)=4-8=-4.
\]

Dus:
\[
\boxed{(-\sqrt2,-4)}
\qquad\text{en}\qquad
\boxed{(\sqrt2,-4)}.
\]

Bij $x=0$ ligt een lokaal maximum:
\[
\boxed{(0,0)}.
\]

De tweede afgeleide is
\[
f''(x)=12x^2-8
=4(3x^2-2).
\]

Dus
\[
f''(x)=0
\quad\Longleftrightarrow\quad
x=\pm\sqrt{\frac23}.
\]

Voor
\[
|x|>\sqrt{\frac23}
\]
is $f''(x)>0$ en is de grafiek hol. Voor
\[
|x|<\sqrt{\frac23}
\]
is $f''(x)<0$ en is de grafiek bol.

De kromming verandert bij beide nulpunten van $f''$. Omdat
\[
f\left(\pm\sqrt{\frac23}\right)
=
\frac49-\frac83
=
-\frac{20}{9},
\]
zijn de buigpunten
\[
\boxed{\left(-\sqrt{\frac23},-\frac{20}{9}\right)}
\]
en
\[
\boxed{\left(\sqrt{\frac23},-\frac{20}{9}\right)}.
\]

De kwalitatieve grafiek is symmetrisch ten opzichte van de $y$-as en moet
alle gevonden extrema en buigpunten respecteren.

\end{oplossing}

\end{oefening}


% ------------------------------------------------------------
% OEFENING 16
% ------------------------------------------------------------

\begin{oefening}[Een redenering beoordelen]{\oefU\ESS}

Een leerling onderzoekt een functie $f$ en vindt:

\[
f'(2)=0.
\]

De leerling besluit:

\begin{quote}
De grafiek van $f$ heeft dus bij $x=2$ een maximum of minimum.
\end{quote}

\begin{question}
Is deze redenering geldig?
\end{question}

\begin{question}
Welke informatie ontbreekt?
\end{question}

\begin{question}
Geef een voorbeeld van een functie waarvoor $f'(2)=0$
en $x=2$ wel bij een extremum hoort.
\end{question}

\begin{question}
Geef een voorbeeld van een functie waarvoor $f'(2)=0$
maar $x=2$ niet bij een extremum hoort.
\end{question}

\begin{oplossing}

De redenering is niet geldig.

Uit
\[
f'(2)=0
\]
volgt alleen dat $x=2$ een stationaire waarde is en dat de raaklijn daar
horizontaal is.

Om een extremum te besluiten, moet onderzocht worden of $f'$ van teken
verandert rond $x=2$.

Een voorbeeld waarbij er wel een extremum is:
\[
f(x)=(x-2)^2.
\]
Dan
\[
f'(x)=2(x-2),
\]
en $f'$ verandert bij $x=2$ van negatief naar positief. Er ligt dus een
lokaal minimum.

Een voorbeeld zonder extremum:
\[
g(x)=(x-2)^3.
\]
Dan
\[
g'(x)=3(x-2)^2.
\]
Deze afgeleide is aan beide kanten positief, zodat de functie blijft
stijgen. Er is dus geen extremum.

\end{oplossing}

\end{oefening}


% ------------------------------------------------------------
% OEFENING 17
% ------------------------------------------------------------

\begin{oefening}[Nog een redenering beoordelen]{\oefU\ESS}

Een leerling vindt voor een functie:

\[
f''(-1)=0.
\]

Daaruit besluit hij dat de grafiek bij $x=-1$ een buigpunt heeft.

\begin{question}
Leg uit waarom deze conclusie niet noodzakelijk correct is.
\end{question}

\begin{question}
Wat moet bijkomend onderzocht worden?
\end{question}

\begin{question}
Geef een voorbeeld waarbij $f''(-1)=0$ en er wel een buigpunt is.
\end{question}

\begin{question}
Geef een voorbeeld waarbij $f''(-1)=0$ maar er geen buigpunt is.
\end{question}

\begin{oplossing}

De conclusie is niet noodzakelijk correct.

Uit
\[
f''(-1)=0
\]
volgt alleen dat $x=-1$ een mogelijke kandidaat voor een buigpunt is.
Er moet bijkomend worden onderzocht of de kromming verandert, bijvoorbeeld
door het teken van $f''$ links en rechts van $-1$ te vergelijken.

Een voorbeeld met een buigpunt is
\[
f(x)=(x+1)^3.
\]
Dan
\[
f''(x)=6(x+1),
\]
zodat $f''$ van negatief naar positief verandert bij $x=-1$.

Een voorbeeld zonder buigpunt is
\[
g(x)=(x+1)^4.
\]
Dan
\[
g''(x)=12(x+1)^2.
\]
Hoewel
\[
g''(-1)=0,
\]
is $g''$ aan beide kanten positief. De kromming verandert dus niet.

\end{oplossing}

\end{oefening}


% ------------------------------------------------------------
% OEFENING 18
% ------------------------------------------------------------

\begin{oefening}[Van tabel naar grafiek]{\oefU\ESS}

Van een functie $f$ kennen we:

\[
\begin{array}{c|ccccccc}
x
& -\infty & & -1 & & 2 & & +\infty \\ \hline

f'(x)
& & + & 0 & - & 0 & + &
\end{array}
\]

en

\[
\begin{array}{c|ccccc}
x
& -\infty & & 0 & & +\infty \\ \hline

f''(x)
& & - & 0 & + &
\end{array}
\]

Verder geldt:

\[
f(-1)=3,\qquad f(0)=1,\qquad f(2)=-2.
\]

\begin{question}
Bepaal de stijg- en daalintervallen.
\end{question}

\begin{question}
Bepaal de lokale extrema.
\end{question}

\begin{question}
Bepaal waar de kromming verandert.
\end{question}

\begin{question}
Bepaal het buigpunt.
\end{question}

\begin{question}
Schets een mogelijke grafiek van $f$ die met alle gegevens
overeenkomt.
\end{question}

\begin{oplossing}

Uit de tekentabel van $f'$ volgt:

\[
f'(x)>0
\]
voor
\[
x<-1
\quad\text{en}\quad
x>2.
\]

Dus $f$ is stijgend op
\[
]-\infty,-1[
\quad\text{en}\quad
]2,+\infty[.
\]

Voor
\[
-1<x<2
\]
geldt
\[
f'(x)<0,
\]
dus $f$ is daar dalend.

Bij $x=-1$ verandert $f'$ van positief naar negatief, dus ligt er een lokaal
maximum. Omdat
\[
f(-1)=3,
\]
is dit
\[
\boxed{(-1,3)}.
\]

Bij $x=2$ verandert $f'$ van negatief naar positief, dus ligt er een lokaal
minimum:
\[
\boxed{(2,-2)}.
\]

Uit de tekentabel van $f''$ volgt dat de kromming verandert bij $x=0$.
Omdat
\[
f(0)=1,
\]
is het buigpunt
\[
\boxed{(0,1)}.
\]

Een mogelijke schets stijgt dus tot $(-1,3)$, daalt vervolgens door het
buigpunt $(0,1)$ tot $(2,-2)$ en stijgt daarna opnieuw.

\end{oplossing}

\end{oefening}


% ============================================================
% TOEPASSINGEN
% ============================================================




% ------------------------------------------------------------
% OEFENING 19
% ------------------------------------------------------------

\begin{oefening}[Beweging langs een rechte]{\oefU\ESS}

De positie van een bewegend voorwerp wordt beschreven door

\[
s(t)=t^3-6t^2+9t,
\qquad t\geq 0,
\]

waarbij $s$ in meter en $t$ in seconde wordt uitgedrukt.

\begin{question}
Bepaal de snelheidsfunctie $v(t)$.
\end{question}

\begin{question}
Op welke tijdstippen staat het voorwerp ogenblikkelijk stil?
\end{question}

\begin{question}
Onderzoek het teken van $v(t)$.
\end{question}

\begin{question}
Wanneer beweegt het voorwerp in de positieve richting?
Wanneer in de negatieve richting?
\end{question}

\begin{question}
Bepaal de versnellingsfunctie $a(t)$.
\end{question}

\begin{question}
Wanneer is de versnelling nul?
\end{question}

\begin{question}
Verbind je resultaten met het verloop van de positiegrafiek.
\end{question}

\begin{oplossing}

Gegeven:
\[
s(t)=t^3-6t^2+9t.
\]

De snelheid is
\[
v(t)=s'(t)=3t^2-12t+9
=3(t-1)(t-3).
\]

Het voorwerp staat ogenblikkelijk stil als
\[
v(t)=0.
\]
Dus:
\[
\boxed{t=1\ \mathrm{s}}
\qquad\text{en}\qquad
\boxed{t=3\ \mathrm{s}}.
\]

Voor $t\geq0$ is het tekenverloop:
\[
\begin{array}{c|ccccccc}
t & 0 & & 1 & & 3 & & +\infty\\ \hline
v(t) & & + & 0 & - & 0 & + &
\end{array}
\]

Het voorwerp beweegt dus in de positieve richting voor
\[
0\leq t<1
\quad\text{en}\quad
t>3,
\]
en in de negatieve richting voor
\[
1<t<3.
\]

De versnelling is
\[
a(t)=v'(t)=6t-12=6(t-2).
\]

Dus
\[
a(t)=0
\quad\Longleftrightarrow\quad
\boxed{t=2\ \mathrm{s}}.
\]

De positiegrafiek stijgt waar $v>0$ en daalt waar $v<0$. Ze heeft een lokaal
maximum bij $t=1$ en een lokaal minimum bij $t=3$.

Verder is
\[
a(t)<0 \quad \text{voor }t<2
\]
en
\[
a(t)>0 \quad \text{voor }t>2.
\]
De kromming van de positiegrafiek verandert dus bij $t=2$.

\end{oplossing}

\end{oefening}


% ------------------------------------------------------------
% OEFENING 20
% ------------------------------------------------------------

\begin{oefening}[Een groeiproces interpreteren]{\oefU}

De hoeveelheid biomassa in een experimenteel systeem wordt tijdens
een bepaalde periode gemodelleerd door

\[
B(t)=-t^3+9t^2+24,
\qquad 0\leq t\leq 6,
\]

waarbij $B$ in gram en $t$ in dagen wordt uitgedrukt.

\begin{question}
Bepaal $B'(t)$.
\end{question}

\begin{question}
Wanneer neemt de biomassa toe en wanneer neemt ze af?
\end{question}

\begin{question}
Bepaal $B''(t)$.
\end{question}

\begin{question}
Wanneer neemt de groeisnelheid toe en wanneer neemt ze af?
\end{question}

\begin{question}
Leg het verschil uit tussen de uitspraken
``de biomassa neemt toe'' en
``de groeisnelheid neemt toe''.
\end{question}

\begin{oplossing}

Gegeven:
\[
B(t)=-t^3+9t^2+24,
\qquad 0\leq t\leq6.
\]

De eerste afgeleide is
\[
B'(t)=-3t^2+18t
=3t(6-t).
\]

Voor
\[
0<t<6
\]
geldt
\[
B'(t)>0.
\]
De biomassa neemt dus gedurende het volledige inwendige interval toe.
Aan de uiteinden geldt
\[
B'(0)=B'(6)=0.
\]

De tweede afgeleide is
\[
B''(t)=-6t+18
=6(3-t).
\]

Voor
\[
0<t<3
\]
geldt
\[
B''(t)>0,
\]
dus de groeisnelheid $B'$ neemt toe.

Voor
\[
3<t<6
\]
geldt
\[
B''(t)<0,
\]
dus de groeisnelheid neemt af.

``De biomassa neemt toe'' betekent
\[
B'(t)>0.
\]

``De groeisnelheid neemt toe'' betekent
\[
B''(t)>0.
\]

De biomassa kan dus nog steeds toenemen terwijl haar groeisnelheid al
afneemt.

\end{oplossing}

\end{oefening}


% ============================================================
% GEVORDERD
% ============================================================




% ------------------------------------------------------------
% OEFENING 21
% ------------------------------------------------------------

\begin{oefening}[Ontwerp een functie]{\oefG}

Construeer een veeltermfunctie $f$ die voldoet aan:

\[
f'(-1)=0,
\qquad
f'(2)=0,
\]

waarbij $f$:

\begin{itemize}
  \item stijgt voor $x<-1$;
  \item daalt voor $-1<x<2$;
  \item opnieuw stijgt voor $x>2$.
\end{itemize}

\begin{question}
Construeer eerst een geschikte formule voor $f'(x)$.
\end{question}

\begin{question}
Bepaal vervolgens een mogelijke functie $f$.
\end{question}

\begin{question}
Controleer met een tekenonderzoek dat je functie aan alle
voorwaarden voldoet.
\end{question}

\begin{question}
Is je antwoord uniek? Leg uit.
\end{question}

\begin{oplossing}

We willen nulpunten van $f'$ bij
\[
x=-1
\qquad\text{en}\qquad
x=2,
\]
met tekenverloop
\[
+\quad -\quad +.
\]

Een eenvoudige keuze is daarom
\[
\boxed{f'(x)=(x+1)(x-2)}.
\]

Inderdaad:
\[
f'(x)>0 \quad \text{voor }x<-1,
\]
\[
f'(x)<0 \quad \text{voor }-1<x<2,
\]
en
\[
f'(x)>0 \quad \text{voor }x>2.
\]

Omdat
\[
f'(x)=x^2-x-2,
\]
is een mogelijke functie
\[
\boxed{
f(x)=\frac13x^3-\frac12x^2-2x
}.
\]

Algemener mag daar nog een constante bij:
\[
f(x)=\frac13x^3-\frac12x^2-2x+C.
\]

Het antwoord is dus niet uniek. Bovendien kan ook $f'$ met een positieve
constante factor vermenigvuldigd worden zonder het gewenste tekenverloop te
veranderen.

\end{oplossing}

\end{oefening}


% ------------------------------------------------------------
% OEFENING 22
% ------------------------------------------------------------

\begin{oefening}[Parameter en stationair punt]{\oefG}

Gegeven is de familie functies

\[
f_a(x)=x^3-3ax.
\]

\begin{question}
Bepaal $f_a'(x)$.
\end{question}

\begin{question}
Voor welke waarden van $a$ heeft de functie twee verschillende
stationaire punten?
\end{question}

\begin{question}
Voor welke waarde van $a$ vallen de stationaire punten samen?
\end{question}

\begin{question}
Onderzoek wat er in dat bijzondere geval met het verloop van de
functie gebeurt.
\end{question}

\begin{oplossing}

Gegeven:
\[
f_a(x)=x^3-3ax.
\]

De afgeleide is
\[
f_a'(x)=3x^2-3a
=3(x^2-a).
\]

Stationaire punten voldoen aan
\[
x^2=a.
\]

Er zijn twee verschillende reële oplossingen
\[
x=\pm\sqrt a
\]
als
\[
\boxed{a>0}.
\]

De twee stationaire punten vallen samen als
\[
\boxed{a=0}.
\]

Dan is
\[
f_0(x)=x^3
\]
en
\[
f_0'(x)=3x^2.
\]

De afgeleide is zowel links als rechts van $0$ positief. De functie blijft
dus stijgen en heeft bij $x=0$ geen extremum. Er is daar een stationair punt
met horizontale raaklijn.

Voor $a<0$ zijn er geen reële stationaire punten.

\end{oplossing}

\end{oefening}


% ------------------------------------------------------------
% OEFENING 23
% ------------------------------------------------------------

\begin{oefening}[Kan zo'n functie bestaan?]{\oefG}

Onderzoek telkens of een functie met de gegeven eigenschappen
kan bestaan. Motiveer grondig.

\begin{question}
Op een interval geldt tegelijk
\[
f'(x)>0
\qquad\text{en}\qquad
f''(x)<0.
\]
\end{question}

\begin{question}
Op een interval geldt
\[
f'(x)<0,
\]
terwijl $f'$ zelf stijgt.
\end{question}

\begin{question}
Een functie heeft in $x=a$ een lokaal minimum, maar
\[
f'(a)\neq 0,
\]
terwijl $f$ in een open interval rond $a$ afleidbaar is.
\end{question}

\begin{question}
Een functie heeft een buigpunt bij $x=a$, maar
\[
f''(a)\neq0.
\]
Onderzoek zorgvuldig welke bijkomende aannames nodig zijn om
hierover te kunnen besluiten.
\end{question}

\begin{oplossing}

\textbf{1.}
\[
f'(x)>0
\qquad\text{en}\qquad
f''(x)<0.
\]

Dit kan. De functie stijgt, maar de hellingen worden kleiner. Bijvoorbeeld
\[
f(x)=-x^2+4x
\]
op het interval $]0,1[$. Daar geldt
\[
f'(x)=4-2x>0
\]
en
\[
f''(x)=-2<0.
\]

\medskip

\textbf{2.}
\[
f'(x)<0
\]
terwijl $f'$ stijgt.

Ook dit kan. Dan is
\[
f''(x)>0.
\]
Bijvoorbeeld
\[
f(x)=x^2
\]
op een interval met $x<0$. Daar is
\[
f'(x)=2x<0
\]
en
\[
f''(x)=2>0.
\]

\medskip

\textbf{3.} Een lokaal minimum met
\[
f'(a)\neq0
\]
terwijl $f$ in een open interval rond $a$ afleidbaar is.

Dit kan niet. Volgens de noodzakelijke voorwaarde van Fermat geldt voor een
afleidbare functie in een inwendig lokaal extremum:
\[
f'(a)=0.
\]

\medskip

\textbf{4.} Een buigpunt bij $x=a$ met
\[
f''(a)\neq0.
\]

Als $f$ tweemaal afleidbaar is in $a$ en een buigpunt betekent dat de
kromming werkelijk verandert, dan kan dit niet: bij zo'n glad buigpunt moet
\[
f''(a)=0.
\]

Een buigpunt kan wel bestaan wanneer $f''(a)$ niet bestaat. Bijvoorbeeld:
\[
f(x)=x|x|.
\]
De kromming verandert bij $x=0$, terwijl de tweede afgeleide daar niet
bestaat.

De bijkomende aannames over het bestaan van $f''(a)$ en de gebruikte
definitie van een buigpunt zijn dus essentieel.

\end{oplossing}

\end{oefening}


% ------------------------------------------------------------
% OEFENING 24
% ------------------------------------------------------------

\begin{oefening}[Reconstrueer het verloop]{\oefG}

Van een functie $f$ is de afgeleide

\[
f'(x)=(x+2)(x-1)^2(x-3).
\]

Je hoeft $f$ zelf niet te bepalen.

\begin{question}
Bepaal alle stationaire waarden van $x$.
\end{question}

\begin{question}
Stel een tekentabel op voor $f'$.
\end{question}

\begin{question}
Bepaal waar $f$ stijgt en daalt.
\end{question}

\begin{question}
Bij welke stationaire punten heeft $f$ een extremum?
\end{question}

\begin{question}
Bij welk stationair punt is er geen extremum?
Leg uit waarom de dubbele factor hierbij belangrijk is.
\end{question}

\begin{question}
Schets een mogelijke kwalitatieve grafiek van $f$.
\end{question}

\begin{oplossing}

Gegeven:
\[
f'(x)=(x+2)(x-1)^2(x-3).
\]

De stationaire waarden zijn
\[
\boxed{x=-2,\qquad x=1,\qquad x=3}.
\]

Omdat
\[
(x-1)^2\geq0,
\]
wordt het teken vooral bepaald door
\[
(x+2)(x-3).
\]

Het tekenverloop is:
\[
\begin{array}{c|ccccccccc}
x & -\infty & & -2 & & 1 & & 3 & & +\infty\\ \hline
f'(x) & & + & 0 & - & 0 & - & 0 & + &
\end{array}
\]

Dus $f$ stijgt op
\[
]-\infty,-2[
\quad\text{en}\quad
]3,+\infty[
\]
en daalt op
\[
]-2,1[
\quad\text{en}\quad
]1,3[.
\]

Bij $x=-2$ verandert $f'$ van positief naar negatief. Daar ligt een lokaal
maximum.

Bij $x=3$ verandert $f'$ van negatief naar positief. Daar ligt een lokaal
minimum.

Bij $x=1$ verandert het teken van $f'$ niet. Er is daar dus wel een
stationair punt, maar geen extremum.

De dubbele factor
\[
(x-1)^2
\]
is precies de reden waarom het teken bij $x=1$ niet verandert.

Een mogelijke kwalitatieve grafiek stijgt tot $x=-2$, daalt vervolgens
door het stationaire punt bij $x=1$ tot $x=3$ en stijgt daarna opnieuw.

\end{oplossing}

\end{oefening}


% ------------------------------------------------------------
% OEFENING 25
% ------------------------------------------------------------

\begin{oefening}[Wat kun je wél besluiten?]{\oefG}

Voor een tweemaal afleidbare functie geldt op een interval:

\[
f'(x)>0
\qquad\text{en}\qquad
f''(x)<0.
\]

Een leerling zegt:

\begin{quote}
De functie stijgt, dus ze stijgt steeds sneller.
\end{quote}

\begin{question}
Leg uit waarom deze uitspraak fout is.
\end{question}

\begin{question}
Beschrijf correct wat er met $f$ gebeurt.
\end{question}

\begin{question}
Beschrijf wat er met $f'$ gebeurt.
\end{question}

\begin{question}
Schets een mogelijke grafiek die aan de voorwaarden voldoet.
\end{question}

\begin{question}
Geef, indien mogelijk, een eenvoudig functievoorschrift dat
op een geschikt interval aan de voorwaarden voldoet.
\end{question}

\begin{oplossing}

De uitspraak van de leerling is fout.

Uit
\[
f'(x)>0
\]
volgt dat de functie stijgt.

Maar uit
\[
f''(x)<0
\]
volgt dat $f'$ afneemt. De positieve hellingen worden dus steeds kleiner.

De functie stijgt bijgevolg
\[
\boxed{\text{steeds minder snel}}.
\]

De grafiek is bol (concaaf).

Een eenvoudig voorbeeld is
\[
f(x)=-x^2+4x
\]
op bijvoorbeeld het interval
\[
]0,1[.
\]

Daar geldt
\[
f'(x)=4-2x>0
\]
en
\[
f''(x)=-2<0.
\]

De grafiek stijgt op dat interval, maar wordt geleidelijk minder steil.

\end{oplossing}

\end{oefening}


% ============================================================
% SYNTHESE
% ============================================================




% ------------------------------------------------------------
% OEFENING 26
% ------------------------------------------------------------

\begin{oefening}[Van formule naar volledige grafiek]{\oefU\ESS}

Gegeven:

\[
f(x)=x^4-4x^3.
\]

Voer een verlooponderzoek uit.

\begin{question}
Bepaal $f'(x)$ en $f''(x)$.
\end{question}

\begin{question}
Bepaal alle stationaire punten.
\end{question}

\begin{question}
Stel een volledige tekentabel op voor $f'$.
\end{question}

\begin{question}
Bepaal de stijg- en daalintervallen.
\end{question}

\begin{question}
Bepaal en classificeer de lokale extrema.
\end{question}

\begin{question}
Onderzoek het teken van $f''$.
\end{question}

\begin{question}
Bepaal de buigpunten.
\end{question}

\begin{question}
Schets de grafiek op basis van je onderzoek.
Markeer de stationaire punten en buigpunten.
\end{question}

\begin{question}
Controleer achteraf je schets met ICT.

Welke informatie uit je algebraïsche onderzoek kun je rechtstreeks
op de grafiek herkennen?
\end{question}

\begin{oplossing}

Gegeven:
\[
f(x)=x^4-4x^3.
\]

De eerste afgeleide is
\[
f'(x)=4x^3-12x^2
=4x^2(x-3).
\]

De tweede afgeleide is
\[
f''(x)=12x^2-24x
=12x(x-2).
\]

De stationaire waarden zijn
\[
x=0
\qquad\text{en}\qquad
x=3.
\]

Het tekenverloop van $f'$:
\[
\begin{array}{c|ccccccc}
x & -\infty & & 0 & & 3 & & +\infty\\ \hline
f'(x) & & - & 0 & - & 0 & + &
\end{array}
\]

Dus $f$ is dalend op
\[
]-\infty,0[
\quad\text{en}\quad
]0,3[
\]
en stijgend op
\[
]3,+\infty[.
\]

Bij $x=0$ verandert $f'$ niet van teken. Er is dus geen extremum, maar wel
een stationair punt.

Bij $x=3$ verandert $f'$ van negatief naar positief. Daar ligt een lokaal
minimum:
\[
f(3)=81-108=-27.
\]

Dus:
\[
\boxed{(3,-27)}.
\]

Voor de tweede afgeleide:
\[
f''(x)>0 \quad \text{voor }x<0,
\]
\[
f''(x)<0 \quad \text{voor }0<x<2,
\]
en
\[
f''(x)>0 \quad \text{voor }x>2.
\]

De kromming verandert dus bij
\[
x=0
\qquad\text{en}\qquad
x=2.
\]

Omdat
\[
f(0)=0
\]
en
\[
f(2)=16-32=-16,
\]
zijn de buigpunten
\[
\boxed{(0,0)}
\qquad\text{en}\qquad
\boxed{(2,-16)}.
\]

Het punt $(0,0)$ is tegelijk een stationair punt en een buigpunt.

De grafiek daalt van links naar rechts tot aan het minimum $(3,-27)$,
waarbij de kromming verandert in $(0,0)$ en $(2,-16)$, en stijgt daarna.

Bij controle met ICT moeten precies deze stationaire punten, het minimum,
de twee buigpunten en de gevonden stijg-, daal- en krommingsintervallen
zichtbaar zijn.

\end{oplossing}

\end{oefening}


% ------------------------------------------------------------
% OEFENING 27
% ------------------------------------------------------------

\begin{oefening}[Kun je alles verbinden?]{\oefU\ESS}

Vul telkens aan zonder eerst een concreet functievoorschrift te zoeken.

\begin{question}
Als $f'(x)>0$, dan is $f$ \ldots
\end{question}

\begin{question}
Als $f'(x)<0$, dan is $f$ \ldots
\end{question}

\begin{question}
Als $f'$ van positief naar negatief verandert bij $x=a$,
dan heeft $f$ daar \ldots
\end{question}

\begin{question}
Als $f'$ van negatief naar positief verandert bij $x=a$,
dan heeft $f$ daar \ldots
\end{question}

\begin{question}
Als $f''(x)>0$, dan \ldots
\end{question}

\begin{question}
Als $f''(x)<0$, dan \ldots
\end{question}

\begin{question}
Voor een buigpunt moet de \ldots\ van de grafiek veranderen.
\end{question}

\begin{question}
Waarom is
\[
f'(a)=0
\]
geen voldoende voorwaarde voor een extremum?
\end{question}

\begin{question}
Waarom is
\[
f''(a)=0
\]
geen voldoende voorwaarde voor een buigpunt?
\end{question}

\begin{oplossing}

Als
\[
f'(x)>0,
\]
dan is $f$ stijgend.

Als
\[
f'(x)<0,
\]
dan is $f$ dalend.

Als $f'$ bij $x=a$ van positief naar negatief verandert, dan heeft $f$
daar een lokaal maximum.

Als $f'$ bij $x=a$ van negatief naar positief verandert, dan heeft $f$
daar een lokaal minimum.

Als
\[
f''(x)>0,
\]
dan is $f'$ stijgend en is de grafiek van $f$ hol (convex).

Als
\[
f''(x)<0,
\]
dan is $f'$ dalend en is de grafiek van $f$ bol (concaaf).

Voor een buigpunt moet de kromming van de grafiek veranderen.

De voorwaarde
\[
f'(a)=0
\]
is niet voldoende voor een extremum, omdat $f'$ mogelijk aan beide kanten
van $a$ hetzelfde teken heeft. Er moet een tekenverandering worden
onderzocht.

Ook
\[
f''(a)=0
\]
is niet voldoende voor een buigpunt. Er moet bijkomend worden nagegaan of
$f''$ van teken verandert, of equivalent of de kromming werkelijk verandert.

\end{oplossing}

\end{oefening}

\FVDendExercises

\end{document}